%% ==========================================================================
%%  Approach 5 --- agent induction, and the balance invariant it would need.
%%
%%  Source: attempted 2026-08-03 at the author's request, starting from Route A
%%  of update_4/PS1_note_next_routes.md.  Scripts: update_5/routeA.py,
%%  routeA_adv.py, routeA_cex.py.
%%
%%  STATUS.  The natural invariant is REFUTED at n=3, m=4 with binary additive
%%  costs.  Agent induction itself is not refuted; what is refuted is the only
%%  assignment-invariant invariant available to it, which is what made the
%%  induction attractive in the first place.
%%
%%  READ THIS BEFORE PROPOSING ANY ARC-LOCAL SUFFICIENT CONDITION.  On the
%%  counterexample every good allocation has an arc at -3 or -2, so no bound of
%%  the form "all arcs >= -1" can certify any of them.
%% ==========================================================================
\section{Approach 5: Agent Induction and the Balance Invariant}
\label{sec:approach5}

Every approach so far inducts on \emph{items}. Inducting on \emph{agents}
instead is attractive for chores because a newcomer starts with an empty bundle,
which is the cheapest bundle there is, so the newcomer is the natural sink and
the orientation problem of \Cref{sec:approach1} does not arise. This section
sets up what such an induction would have to carry, and shows that the natural
candidate fails.

% --------------------------------------------------------------------------
\subsection{What the induction would have to carry}

Corollary~\ref{cor:onestep} says that if $\alloc$ is envy-freeable and every arc
satisfies $\edgew{\alloc}(i,j) \ge -1$ --- equivalently
$\cost_i(\bundle{j}) \le \cost_i(\bundle{i}) + 1$ --- then
$\optsubsidy \in \set{0,1}^n$. The useful feature of that condition for an
induction on agents is that it is \emph{assignment-invariant}: it constrains
only the family of bundles, not which agent holds which. That suggests carrying
the family rather than the allocation.

\begin{definition}[Uniformly balanced family]
\label{def:unifbal}
A family of $n$ disjoint bundles covering $\items$ is \emph{uniformly balanced}
if every agent values all $n$ bundles within one unit of each other:
$\max_t \cost_i(B_t) - \min_t \cost_i(B_t) \le 1$ for every $i \in \agents$.
\end{definition}

\begin{proposition}
\label{prop:balance-suffices}
If an instance admits a uniformly balanced family, then
Conjecture~\ref{conj:main} holds for it.
\end{proposition}

\begin{proof}
Assign the bundles to agents by a maximum-weight matching. By
Theorem~\ref{thm:hs-characterisation} the resulting allocation is
envy-freeable, and by uniform balance every arc satisfies
$\edgew{\alloc}(i,j) = \cost_i(\bundle{i}) - \cost_i(\bundle{j}) \ge -1$
--- the bound holds for \emph{any} assignment, so in particular for this one.
Corollary~\ref{cor:onestep} then gives $\optsubsidy \in \set{0,1}^n$, and
Observation~\ref{obs:integrality} plus the vanishing of $\pathw{}$ at the
endpoint of a heaviest path give a total of at most $n-1$.
\end{proof}

This is exactly the shape an agent induction wants: a purely combinatorial,
schedule-free, coverage-free statement, with the step from $n$ to $n+1$ agents
being a refinement of the family. It also holds on every hard instance the
project has accumulated, and on the full exhaustive family at $n = m = 3$ ---
all $9{,}880$ instances, with the implication chain of
Proposition~\ref{prop:balance-suffices} verified end to end on each
(\texttt{routeA.py}). It is nevertheless false.

% --------------------------------------------------------------------------
\subsection{The invariant is unreachable}

\begin{proposition}
\label{prop:no-balance}
There is a negative dichotomous instance with $n = 3$ and $m = 4$, with binary
additive costs, admitting no uniformly balanced family.
\end{proposition}

\begin{proof}
Take $\cost_i(S) = \abs{S \cap D_i}$ with
\[
  D_1 = \set{g_1,g_2}, \qquad D_2 = \set{g_1,g_3,g_4}, \qquad
  D_3 = \set{g_2,g_3,g_4}.
\]
Uniform balance for agent $i$ means the elements of $D_i$ are spread across the
three bundles with counts of range at most $1$. Three elements over three
bundles have range at most $1$ only in the pattern $(1,1,1)$, so $D_2$ forces
$g_1,g_3,g_4$ into three \emph{different} bundles, one each, and $D_3$ likewise
forces $g_2,g_3,g_4$ into three different bundles. Hence $g_2$ must lie in the
one bundle containing neither $g_3$ nor $g_4$, which is the bundle containing
$g_1$. But two elements over three bundles have range at most $1$ only in the
pattern $(1,1,0)$, so $D_1$ forces $g_1$ and $g_2$ into \emph{different}
bundles. Contradiction.
\end{proof}

The instance is binary additive and therefore already settled
by~\cite{LMS26}; what it refutes is the method, on a class the literature
closes. Indeed $19$ of its $81$ allocations are good, several with
$\pathw{} = (0,0,0)$ --- it is an easy instance for the conjecture and an
impossible one for the invariant.

The failure is sharper than mere unreachability.

\begin{proposition}
\label{prop:arc-bound-blind}
On the instance of Proposition~\ref{prop:no-balance}, \emph{every} good
allocation has an arc of weight $-2$ or less. Hence no sufficient condition of
the form ``all arcs $\ge -1$'' can certify any of them.
\end{proposition}

\begin{proof}
Exhaustive over all $3^4 = 81$ allocations (\texttt{routeA\_cex.py}): $19$ are
good, and the minimum arc weight over them is $-3$, attained at
$(\set{g_1,g_3,g_4},\set{g_2},\emptyset)$, with every other good allocation
reaching $-2$.
\end{proof}

\begin{remark}[What this rules out]
\label{rem:arc-vs-path}
Corollary~\ref{cor:onestep} is a bound on \emph{arcs}; the target is a bound on
\emph{paths}. Propositions~\ref{prop:no-balance} and~\ref{prop:arc-bound-blind}
show the gap between the two is not slack that a cleverer construction can close
--- it is where all the surviving witnesses live. \textbf{Any correct invariant
must be path-based.} That is precisely what makes agent induction hard, because
path weights depend on which agent holds which bundle, so the invariant cannot
be assignment-invariant and the newcomer's arrival re-opens the whole graph.
Agent induction is not refuted here; the only invariant that would have made it
easy is.
\end{remark}

\begin{remark}[The obstruction is hypergraph discrepancy]
\label{rem:discrepancy}
For binary additive costs, uniform balance says every $D_i$ is split evenly
across the $n$ bundles, which is a discrepancy condition on the set system
$\set{D_1,\dots,D_n}$. Proposition~\ref{prop:no-balance} is a small
discrepancy obstruction, and it is the same combinatorial phenomenon that drives
the reduction by which~\cite{BSV21} proves deciding exact envy-freeness for
binary additive chores NP-complete. The certified-hard set-splitting instance
recorded from that reduction also admits no uniformly balanced family, checked
by \texttt{routeA\_adv.py}. So the balance signal of Remark~\ref{rem:balance},
which three independent routes had pointed at, cannot be imposed exactly ---
only approximately, and the approximation is exactly the path slack.
\end{remark}

% --------------------------------------------------------------------------
\subsection{What survives: the two-tier restatement}

One reformulation does survive, and it is worth recording because it is
equivalent rather than merely sufficient.

\begin{proposition}
\label{prop:twotier-equiv}
Let $\alloc$ be envy-freeable with every arc at most $1$. Then $\alloc$ is good
if and only if the digraph of weight-$1$ arcs contains no directed path of
length two.
\end{proposition}

\begin{proof}
If no such path exists then every directed path carries at most one positive
arc, so has weight at most $1$, and $\pathw{\alloc} \le 1$. Conversely if
$\pathw{\alloc} \le 1$ then $\optsubsidy \in \set{0,1}^n$ by
Observation~\ref{obs:integrality}; writing $S = \set{i : \optsubsidy_i = 1}$,
Proposition~\ref{prop:twotier} puts every arc within $S$, within
$\agents \setminus S$, or from $\agents \setminus S$ into $S$ at weight at most
$0$, so weight-$1$ arcs run only from $S$ to $\agents \setminus S$ and cannot be
composed.
\end{proof}

So the target is exactly: \emph{find an envy-freeable allocation together with a
bipartition of the agents such that all strict envy runs from one side to the
other}. This is Proposition~\ref{prop:twotier} again, now as a criterion rather
than a certificate format, and it is the one direction in
\Cref{sec:conclusion} still untried --- choose the paid set first, then
allocate. It is not assignment-invariant, consistent with
Remark~\ref{rem:arc-vs-path}, but it is a condition on a bipartition rather than
on a schedule, so it inherits none of the dead ends of
\Cref{sec:approach3} and none of the ex-ante failures of
\Cref{sec:approach4}.
